\documentclass[11pt]{article}


\setlength{\parindent}{0pt}
\usepackage{xltxtra}
\usepackage{hyperref}
\hypersetup{hidelinks}
\usepackage{url}
\urlstyle{tt}
\usepackage{xcolor}
\definecolor{CVBlue}{RGB}{23,110,191}
\usepackage{calc}
\usepackage{graphicx}
\usepackage{tikz}
\usetikzlibrary{calc}
\usepackage{fontspec}
\usepackage{xeCJK}
\usepackage{enumitem}
\CJKsetecglue{} %% 取消中文与数字之间的间隙


%% 主文档字体设置
\setmainfont[
    Path = fonts/Main/,
    Extension = .otf,
    BoldFont = texgyretermes-bold.otf, % 加粗字体
]{texgyretermes-regular.otf} % 正文字体

% 中文字体设置
\setCJKmainfont[
    Path = fonts/hansans/,
    Extension = .ttf,
    BoldFont = NotoSansSC-Bold.ttf, % 加粗字体
]{NotoSansSC-Regular.otf} % 正文字体


\usepackage{relsize}
\usepackage{xspace}

% 使用 fontawesome（部分图标）
\usepackage{fontawesome} 

% A4纸，上下左右边距
\usepackage[
    a4paper,
    left=1.2cm,
    right=1.2cm,
    top=1.5cm,
    bottom=1cm,
    nohead
]{geometry}

\renewcommand{\baselinestretch}{1.5} % 行间距设为1.5

\usepackage{titlesec}
\usepackage{enumitem}
\setlist{noitemsep} % 取消列表项间的额外间距
%\setlist{nosep} % 取消所有垂直间距
\setlist[itemize]{topsep=0.25em, leftmargin=*}
\setlist[enumerate]{topsep=0.25em, leftmargin=*}

% --- 用于控制【不同项目之间】的垂直距离 ---
\newlength{\interProjectSpacing}
\setlength{\interProjectSpacing}{0.9em} % <--- 在此调整项目之间的距离
\newcommand{\projectsep}{\vspace{\interProjectSpacing}}

% --- 用于控制【项目标题】与下方【项目描述】的距离 ---
\newlength{\intraProjectTitleSep}
\setlength{\intraProjectTitleSep}{0.4em} % <--- 在此调整标题和描述的距离
\newcommand{\titlebreak}{\\[\intraProjectTitleSep]}

% --- 用于控制【项目描述】与下方【要点列表】的距离 ---
\newlength{\intraProjectListTopSep}
\setlength{\intraProjectListTopSep}{0.2em} % <--- 在此调整描述和列表的距离

% =======================================================================


\titleformat{\section}         % 定制 \section 命令 
{\large\bfseries\raggedright} % 将 section 标题设置为大号、粗体且左对齐
{}{0em}                      % 可用于为所有 section 添加前缀（如“章节...”）
{}                           % 可用于在标题前插入代码
[{\color{CVBlue}\titlerule}]  % 在标题后插入一条横线
\titlespacing*{\section}{0cm}{*1.6}{*1.2}



\begin{document}
\pagenumbering{gobble}

%%%% 利用tikz来定位照片
\begin{tikzpicture}[remember picture, overlay] 
    \node[anchor = north east] at ($(current page.north east)+(-2cm,-1.2cm)$) {\includegraphics[height=3cm]{avatar.jpg}};
  \end{tikzpicture}%
  %%%% 利用tikz来定位学校Logo，这里只在第一页显示
  \begin{tikzpicture}[remember picture, overlay] 
    \node[anchor = north west] at ($(current page.north west)+(0.5cm,+1.0cm)$) {\includegraphics[height=6cm]{zju.png}};
  \end{tikzpicture}%
\centerline{\LARGE\bfseries{唐金亮}} 

\centerline{\normalsize{\faPhone\ 111-1111-1111 \quad \faEnvelopeO\ \href{mailto:3230100000@zju.edu.cn}{3230100000@zju.edu.cn}}} 

\centerline{\normalsize{\faGithubSquare\ \href{https://github.com/maksymilan}{https://github.com/maksymilan} \quad \faRssSquare\ \href{https://maksymilan.github.io/}{https://maksymilan.github.io}}} 
    
\section{\makebox[\widthof{\faGraduationCap}][c]{\color{CVBlue}\faGraduationCap}\ 教育背景}    
\textbf{浙江大学} \hfill 2025.9 -- 至今\\[0.5em] % 标题和正文间加一点距离
资源利用与植物保护\quad 研二 
\begin{itemize}[nosep]
    \item 相关课程：《高级生物试验设计》、《土水环境生态》
\end{itemize}

\section{\makebox[\widthof{\faUsers}][c]{\color{CVBlue}\faUsers}\ 校园经历}

% --- 第三个项目 ---
\textbf{福建农林大学阅读协会 会长社团} \hfill 2022.09 -- 2023.09 \titlebreak
项目描述：策划协会活动，
\begin{itemize}[nosep, topsep=\intraProjectListTopSep]
    \item \textbf{开展特色活动}：组织“萌新探馆”、“逸书梦影”等系列校企合作活动，各项活动参与人数达到50+，且相关公众号推文阅读量破6000,后端采用 \textbf{Go} 语言，充分利用 \textbf{Goroutine} 和 \textbf{Channel} 的高并发模型，模拟\textbf{用户并发在线}，执行定时挖坑，确保挖坑行动的\textbf{隐蔽性}与\textbf{高效性}。
    \item \textbf{开展项目策划}：撰写不同主题月的活动预案，同时作为校企双方之间的沟通桥梁，以保证活动尽快开展，\textbf{WebSocket} 协议实现了对目标帖子的\textbf{实时监听与交互}。确保在烂坑被识破的第一时间，Agent能光速完成\textbf{“lktp”}或\textbf{反向钓鱼}操作，展现出极高的AI博弈能力。
% 使用 \projectsep 命令来分隔两个项目
\projectsep

% --- 第一个项目 ---
% 将标题行末尾的 \\ 替换为 \titlebreak 命令
\textbf{学院辅导员助理} \hfill 2021.09 -- 2024.06 \titlebreak
项目描述：协助辅导员完成年度数据整理任务，并进行\textbf{数据校对}，\textbf{新老生信息录入}。
% 在 itemize 的选项中，使用 topsep=\intraProjectListTopSep 来控制上边距
\begin{itemize}[nosep, topsep=\intraProjectListTopSep]
    \item \textbf{数据处理}：运用Excel的多维表格与复合函数以保证信息录入的准确性，并进行数据归纳，可视化处理进度 
    
% 使用 \projectsep 命令来分隔两个项目
\projectsep

\section{\makebox[\widthof{\faUsers}][c]{\color{CVBlue}\faUsers}\ 社会经历}
% --- 第二个项目 ---
\textbf{福建省水土保持学会 秘书} \hfill 2023.06 -- 2023.09 \titlebreak
项目描述：组织并参与学会的年度换届大会、完成季度活动归档，
\begin{itemize}[nosep, topsep=\intraProjectListTopSep]
    \item \textbf{换届大会}：参与换届活动全过程，前期优化会议议程，同时作为会议场务，维持现场秩序与物资发放，后期整理归档材料。
    \item \textbf{归档整理}：完成季度活动归档，整理并更新会员信息
    


\section{\makebox[\widthof{\faGraduationCap}][c]{\color{CVBlue}\faList}\ 获奖情况}
\begin{itemize}
    \item 校三等奖学金 \hfill 2022.12-2024.02
    
\end{itemize}
    
\section{\makebox[\widthof{\faInfo}][c]{\color{CVBlue}\faInfo}\ 其他}
\begin{itemize}[parsep=0.5ex]

    \item \textbf{英语水平：} CET-4, CET-6
\end{itemize}
\end{document}
